\section{结论}

本文围绕纳米硒（SeNPs）细胞毒性预测任务，构建了由基准数据集模型比较、纳米硒专用建模以及增强特征自动筛选组成的研究流程。结果表明，在大规模异质纳米颗粒表格数据上，深度表格模型与树模型均能够完成有效预测；而在样本规模有限、材料类型相对固定的纳米硒场景中，树模型表现出更高的稳定性和适用性。

纳米硒专用数据集包含从 43 篇发表于 2008--2022 年的同行评议文献整理出的 383 条有效记录。基础理化特征能够提供初步预测能力，而元素 Se 含量、长期稳定性和 Se 3d 结合能等材料化学增强变量可进一步改善模型拟合能力。本研究采用的多模态融合框架通过基础与增强描述符拼接、交叉注意力加权以及残差连接整合异质特征，在 IC50 等辅助信息可用时动态突出相关描述符。

在所比较的模型中，XGBoost 在增强特征配置下取得最高回归表现。但该配置包含与响应终点相关的辅助变量 IC50，而非仅包含预测前可获得的先验理化描述符，因此其结果更适合作为机理分析与验证的辅助基准，不宜作为前瞻预测的实际性能指标。随机森林则在精度、稳健性和可解释性之间保持了较好平衡，是后续平台部署与工程实现的更合适候选。

特征筛选表明，表现较优的增强变量组合无需包含全部候选描述符；仅使用 IC50、元素 Se 含量、长期稳定性和 Se 3d 结合能即可达到接近完整增强特征集的精度。这提示应区分两类增强变量：材料本征化学描述符，包括元素 Se 含量、稳定性和结合能，在单一材料体系中具有可迁移、可积累的特征工程价值；IC50 等响应相关辅助变量则在历史实验数据可用时提供诊断参考。因此，后续数据采集与实验设计应优先测量能够用于前瞻筛选的材料本征化学属性。

总体而言，本研究形成了由数据整理、模型比较、特征筛选和多模态融合组成的纳米硒毒性预测研究路径。\textbf{XGBoost 增强特征模型}提供了较高的诊断参考指标（$R^2=0.957$），而材料化学增强随机森林更适合作为后续部署的基线。目前正在表征的自合成纳米硒数据将用于后续外部验证，以评价 IC50 不可用的前瞻场景下模型的泛化能力。这些结果为计算辅助纳米毒性评价和实验优先级筛选提供了基础。
